\documentclass[10pt,a4paper]{amsart}
\usepackage[margin=19mm]{geometry}
\usepackage{amsmath,amssymb,mathtools}
\usepackage[scr=rsfs,cal=euler]{mathalfa}
\usepackage{xfrac}
\usepackage[unicode,hidelinks,bookmarks=false]{hyperref}
\newcommand{\ie}{i.e.\ }
\newcommand{\cf}{cf.\ }
\newcommand{\resp}{respectively\ }
\newcommand{\Subsection}{\subsection}
\providecommand{\nobreakdash}{\nobreak\hbox{-}}
\setlength{\emergencystretch}{2em}
\multlinegap0pt
\usepackage{bm}
\usepackage{bbm}
\usepackage{dsfont}
\usepackage{xspace}
\usepackage{cancel}
\usepackage{mathtools}
\usepackage{cleveref}
\usepackage{amsthm}
\makeatletter
\@ifundefined{theorem}{\newtheorem{theorem}{Theorem}}{}
\@ifundefined{proposition}{\newtheorem{proposition}[theorem]{Proposition}}{}
\makeatother
\providecommand{\subs}{\triangleleft}
\title[Boolean degree one functions on the Grassmann scheme]{Boolean degree one functions on the Grassmann scheme}
\author{Yuval Filmus}
\date{Source: arXiv:2606.23465v2 (CC BY 4.0)}
\begin{document}
\maketitle

\noindent\emph{Source and licence.} Boolean degree one functions on the Grassmann scheme. Version arXiv:2606.23465v2 (CC BY 4.0), distributed under the Creative Commons Attribution 4.0 International licence, as recorded in the arXiv licence metadata for this exact version and shown on the abstract page https://arxiv.org/abs/2606.23465v2. Full rights notice: licenses/arxiv-2606.23465-CC-BY-4.0.txt.\par\smallskip

\noindent\textit{Editorial scope.} Counted unit: the Proposition whose source label is \texttt{pro:step3} in the article (source0.tex lines 380--382, source label \texttt{pro:step3}), delivered together with the article's own complete original proof of it (source0.tex lines 419--488, one \texttt{proof} environment of 70 lines). No mathematics is written, repaired or re-derived: statement and proof are the article's own text line for line. The proof is detached from the statement in the source: the article states the result at lines 380--382 and prints its proof at lines 419--488, and the proof environment's own header names the statement, so the association is the article's own. Required context retained uncounted: none. Every definition, display and label of those lines is retained unchanged. Omitted from this package and not redistributed: 1--379, 383--418, 489--699. Nothing inside a retained excerpt is omitted or altered: every retained body is delivered exactly as the article's lines are written, so no omission receipt of any kind accompanies this package. Citations: the retained excerpts carry no citation command at all, so no bibliography entry of the article is transcribed and no reference is redistributed. Reference adaptation: none was needed and none was made; every reference inside the delivered unit resolves inside it, so no unresolved reference or citation is produced. Printed slips kept literal: no printed word, symbol, hypothesis or number of the counted statement or of its proof was altered. Layout and class: the article's own class is \texttt{article} with packages including fullpage, amsmath, amsfonts, amssymb, amsthm, thmtools, hyperref, cleveref, enumerate, color, tcolorbox; the wrapper is the lane's pinned \texttt{amsart} editorial wrapper at 10pt on A4, which restates the theorem-like environments the retained text needs and adds the article's own macro definitions that the retained text needs (\texttt{\textbackslash subs}), with extra wrapper packages (bm, bbm, dsfont, xspace, cancel, mathtools, cleveref). The delivered numbering is the unit's own. Assets: no retained span contains a figure environment, a graphics-inclusion command, a PDF-page inclusion, a table or a TikZ picture; no image file is copied into this package, the package declares no asset dependency and no image file is redistributed with it. Licence: version arXiv:2606.23465v2 of this article is distributed under the Creative Commons Attribution 4.0 International licence, as recorded in the arXiv licence metadata for this exact version and shown on the abstract page https://arxiv.org/abs/2606.23465v2. A full rights notice is delivered at licenses/arxiv-2606.23465-CC-BY-4.0.txt. Characters: every retained excerpt is pure ASCII, so the delivered engine, XeLaTeX with its default Latin Modern text font, reports no missing character.
\begin{proposition} \label{pro:step3}
If $J_q(a;b)$ is DOT and $a,b \ge 2$ then $J_q(a';b)$ is DOT for all $a' \ge a$.    
\end{proposition}

\begin{proof}[Proof of \Cref{pro:step3}]
Let $n = a+b$ and $n' = a'+b$; we can assume that $n' > n$. Let $f$ be a Boolean degree one function on $J_q(n',b)$. For every $n$-dimensional subspace $V$, the restriction $f|_V$ is a Boolean degree one function on $J_q(n,b)$, and so by assumption, $f|_V$ is trivial.

The proof will involve the Grassmann graph $G = G_q(n',n)$ on the vertex set $J_q(n',n)$, where we connect two subspaces $V,V'$ if their intersection has dimension $n-1$. It is well-known that this graph is connected. If $V,V'$ are neighbors we denote this by $V \sim V'$.
In this case $V \cap V'$ is a codimension~$1$ subspace of $V$, which we denote by $V \cap V' \subs V$.

The subgraph $G_{r^-}$ of $G$ consisting of all subspaces not orthogonal to $r$ is also connected (when $n' > n$). To see this, let $V,V' \in G_q(n',n)$ be two subspaces not orthogonal to $r$. Let $p \in V$ and $p' \in V'$ be arbitrary vectors which are not orthogonal to $r$. Since $n \ge 2$, we can find $W \in G_q(n',n)$ containing both $p$ and $p'$. The subgraph $G_p$ of $G$ consisting of all subspaces containing $p$ is a Grassmann graph and so it is connected. Therefore we can find a path in $G_p$ from $V$ to $W$. We can similarly find a path in $G_{p'}$ from $W$ to $V'$. All vertices on the corresponding path from $V$ to $V'$ are subspaces containing either $p$ or $p'$, and so are not orthogonal to $r$.

\smallskip

The first step in the proof involves constructing the ``restriction table'', which describes the possible structures of $f|_U$ given $f|_V$, where $U \subs V \in J_q(n',n)$:
\[
\begin{array}{c|lll}
f|_V & f|_U \\\hline
0^\pm & 0^\pm \\
x_p^\pm & 0^\pm \text{ if } p \notin U, & x_p^\pm \text{ if } p \in U \\
y_r^\pm & 0^\mp \text{ if } r \perp U, & y_r^\pm \text{ if } r \not\perp U \\
(x_p+y_r)^\pm & 0^\mp \text{ if } r \perp U, & (x_p+y_r)^\pm \text{ if } p \in U, & y_r^\pm \text{ otherwise} 
\end{array}
\]

% The neighbor table describes the possible structured of $f|_{V'}$ given $f|_V$, where $V' \sim V$:
% \[
% \begin{array}{c|l}
% f|_V & f|_{V'} \\\hline
% 0^\pm & 0^\pm, x_p^\pm, y_r^\mp, (x_p+y_r)^\mp\\
% x_p^\pm & 0^\pm, x_{p'}^\pm, y_r^\mp, (x_{p'}+y_r)^\mp \\
% y_r^\pm & 0^\mp, x_p^\mp, y_{r'}^\pm, (x_p+y_{r'})^\pm \\
% (x_p+y_r)^\pm &
% \end{array}
% \]

Let us say that a trivial function is of \emph{positive type} if it is of the form $0,x_p,y_r^-,(x_p+y_r)^-$. Otherwise, it is of \emph{negative type}. If $V \sim V'$ then $V \cap V' \subs V,V'$, and so the restriction table implies that $f|_V$ and $f|_{V'}$ have the same type. Since the Grassmann graph is connected, we see that either all $f|_V$ are of positive type, or all $f|_V$ are of negative type. Without loss of generality, assume that all $f|_V$ are of positive type.

In order to cut the case analysis,\footnote{We thank ChatGPT~5.6~Sol for this idea.} we write each function of positive type as $\alpha(f|_V) + \beta(f|_V)$, where $\alpha(f|_V)$ is either $0$ or of the form $\pm x_p$, and $\beta(f|_V)$ is either $0$ or of the form $y_r^-$. The corresponding restriction tables (derived from the previous restriction table) are

\[
\begin{array}[t]{c|l}
\alpha(f|_V) & \alpha(f|_U) \\\hline
0 & 0 \\
\pm x_p & 0 \text{ if } p \in U, \\ & \pm x_p \text{ if } p \notin U
\end{array}
\quad
\begin{array}[t]{c|l}
\beta(f|_V) & \beta(f|_U) \\\hline
0 & 0 \\
y_r^- & 0 \text{ if } r \perp U, \\ & y_r^- \text{ if } r \not\perp U
\end{array}
\]

We determine the possible behaviors of $\alpha(f|_V)$ and $\beta(f|_V)$ separately.

If $\alpha(f|_{V_0}) = \pm x_p$ for some $V_0,p$ then $\alpha(f|_V) = \pm x_p$ whenever $p \in V$ (with the same sign).
Indeed, consider the subgraph $G_p$ of $G$ consisting of all subspaces containing $p$, which is isomorphic to $G_q(n'-1,n-1)$, and is therefore connected. The restriction table implies that for $V \stackrel{G_p}\sim V'$, if $\alpha(f|_V) = \pm x_p$ then $\alpha(f|_{V'}) = \pm x_p$ as well. Since $\alpha(f|_{V_0}) = \pm x_p$ and $G_p$ is connected, it follows that $\alpha(f|_V) = \pm x_p$ for all $V \in G_p$.

We claim that moreover, $\alpha(f|_V) = 0$ whenever $p \notin V$. To see this, suppose that $\alpha(f|_V) \in \{ x_{p'}, -x_{p'} \}$. Since $n \ge 2$, we can find $W \in J_q(n',n)$ containing both $p$ and $p'$; but then $\alpha(f|_W) = \pm x_p$ and $\alpha(f|_W) \in \{x_{p'}, -x_{p'}\}$, and we reach a contradiction.

Similarly, if $\beta(f|_{V_0}) = y_r^-$ for some $V_0,r$ then $\beta(f|_V) = y_r^-$ whenever $r \not\perp V$.
Indeed, consider the subgraph $G_{r^-}$ of $G$ consisting of all subspaces not perpendicular to $r$, which is connected as shown above. The restriction table implies that for $V \stackrel{G_{r^-}}\sim V'$, if $\beta(f|_V) = y_r^-$ then $\beta(f|_{V'}) = y_r^-$ as well. Since $\beta(f|_{V_0}) = y_r^-$ and $G_{r^-}$ is connected, it follows that $\beta(f|_V) = y_r^-$ for all $V \in G_{r^-}$.

We claim that moreover, $\beta(f|_V) = 0$ whenever $r \perp V$. To see this, suppose that $\beta(f|_V) = y_{r'}^-$. Since $n \ge 2$, we can find a point $p \not\perp r,r'$ and a subspace $W \in J_q(n',n)$ containing $p$; but then $\beta(f|_W)$ would be equal to both $y_r^-$ and $y_{r'}^-$, and we reach a contradiction.

Considering the possible ways in which $\alpha(f|_V)$ and $\beta(f|_V)$ combine to make $f|_V$, we see that one of the following must hold:
\begin{itemize}
\item $\alpha(f|_V) = \beta(f|_V) = 0$ for all $V$. In this case, $f = 0$.
\item $\alpha(f|_V) \in \{0, x_p\}$, according to whether $p \in V$, and $\beta(f|_V) = 0$ for all $V$. In this case, $f = x_p$.
\item $\alpha(f|_V) = 0$ for all $V$, and $\beta(f|_V) \in \{0, y_r^-\}$, according to whether $r \perp V$. In this case, $f = y_r^-$.
\item $\alpha(f|_V) \in \{0, -x_p\}$, according to whether $p \in V$ or not, and $\beta(f|_V) \in \{0, y_r^-\}$, according to whether $r \perp V$ or not. Since $r \perp V$ must imply $p \notin V$, necessarily $p \not\perp r$. In this case, $f = (x_p + y_r)^-$. \qedhere
\end{itemize}
\end{proof}

\end{document}
